\documentclass[journal]{IEEEtran}
\usepackage{graphicx}
\usepackage{amsmath}
\usepackage{listings}
\usepackage{xcolor}
\usepackage{hyperref}
\usepackage{booktabs}

\lstset{
  language=Python,
  basicstyle=\ttfamily\small,
  keywordstyle=\color{blue},
  commentstyle=\color{gray},
  stringstyle=\color{orange},
  breaklines=true,
  frame=single,
  numbers=left,
  numberstyle=\tiny\color{gray}
}

\title{Medial Axis Detection of a Moving Surgical Tool\\Using a Custom Hough Line Transform}
\author{Neal Kauntia\\ ARK Fresher Selection Task 2.3}

\begin{document}
\maketitle

\begin{abstract}
This document describes a complete pipeline for detecting and visualizing the medial 
axis (central skeletal line) of a moving surgical tool in a video sequence. The 
approach consists of four stages: background subtraction using MOG2, morphological 
image cleaning, Sobel-based edge detection, and a from-scratch Hough Line Transform 
implementation. The two strongest parallel edge lines are detected in the Hough 
accumulator, and their geometric midpoint defines the medial axis, which is overlaid 
on the original video frames. The algorithm was developed and debugged without using 
any built-in Hough transform functions, in compliance with the task constraints.
\end{abstract}

\section{Introduction}
Detecting the medial axis of a thin, elongated surgical tool in video is a problem 
with direct applications in robotic surgery, augmented reality overlays, and 
instrument tracking. The tool appears as a narrow, bright, approximately straight 
object moving against a cluttered background. Its medial axis — the line equidistant 
from both long edges — defines the tool's orientation and position.

The core algorithmic challenge is computing this axis without any built-in line 
detection functions. This requires implementing the Hough Line Transform from 
mathematical first principles: constructing the $(\rho, \theta)$ accumulator, 
voting for each edge pixel, and interpreting the accumulator to extract line parameters.

\section{Problem Statement}
Given a video of a moving surgical tool:
\begin{itemize}
    \item \textbf{Input:} Raw RGB video frames containing the surgical tool against 
    a mostly static background.
    \item \textbf{Objective:} For each frame, detect the medial axis of the tool 
    and draw it as a green line overlay on the original frame.
    \item \textbf{Pipeline:} Background subtraction $\rightarrow$ image cleaning 
    $\rightarrow$ edge detection $\rightarrow$ Hough transform $\rightarrow$ 
    medial axis computation.
    \item \textbf{Constraint:} \texttt{cv.HoughLines} and \texttt{cv.HoughLinesP} 
    are forbidden. The full Hough accumulator must be built manually.
\end{itemize}

\section{Related Work}
\textbf{MOG2 Background Subtraction} models each pixel as 
a mixture of Gaussians, adaptively updating the model over time. It robustly 
separates foreground objects from a dynamic background and is the standard method 
for surgical instrument segmentation in video.

\textbf{Hough Line Transform} maps each edge pixel 
$(x, y)$ to a sinusoidal curve in $(\rho, \theta)$ parameter space via the 
relation $\rho = x\cos\theta + y\sin\theta$. Lines in image space correspond 
to peaks in the voting array. The standard form uses $\rho \in [-d, d]$ where 
$d$ is the image diagonal, and $\theta \in [0°, 180°)$.

\textbf{Medial Axis} of a shape is the locus of centres of 
maximal inscribed circles. For a thin elongated shape such as a surgical tool, 
this reduces to the centreline between the two parallel long edges, which can 
be computed as the average of the two detected parallel Hough lines.

\section{Initial Attempts}

\subsection{Attempt 1: Single Peak Extraction}
The first version of the Hough analysis extracted only the single strongest 
peak from the vote

and used that as the medial axis directly.

\begin{lstlisting}
best_index = np.argmax(vote.flatten())
rho_index, theta_index = np.unravel_index(best_index, vote.shape)
\end{lstlisting}

\textbf{Problem:} The strongest peak corresponds to one edge of the tool, 
not the centre. The resulting line was consistently offset to one side of 
the tool. No medial axis was computed — this was simply the most prominent edge.

\subsection{Attempt 2: Column-Sum Based Theta Selection with Rho Averaging}
The second version searched for the dominant $\theta$ by summing vote 
columns (all $\rho$ values for each $\theta$), then extracted two rho values 
at that theta.

\textbf{Bug discovered:} The selected \texttt{medialtheta} was stored as a 
raw index (integer 0--179) rather than the actual angle in radians. Passing 
this integer to \texttt{numpy.cos()} computed the cosine of 0--179 \emph{radians} 
(wrapping many times around the unit circle), producing a completely wrong 
line direction. This was the primary cause of the line appearing perpendicular 
to the tool. The fix was to convert: \texttt{best\_theta = thetas[best\_theta\_index]}.

A secondary bug: the along-line direction vector was computed as 
\texttt{dx = -b, dy = a} with \texttt{a = cos(theta), b = sin(theta)}. 
This gives the \emph{normal} direction (perpendicular to the line), not the 
tangent. The correct tangent direction is \texttt{dx = -sin(theta), dy = cos(theta)}.

\subsection{Attempt 3: Parallel Line Search for True Medial Axis}
After fixing the above bugs, the approach was extended to search for two 
parallel lines (similar $\theta$, different $\rho$) to properly define 
the medial axis as their average.

\begin{lstlisting}
# Sorted vote peaks
index = np.argsort(vote.flatten())[::-1]
rho_index1, theta_index1 = np.unravel_index(index[0], vote.shape)

# Finding second parallel line
for i in indices[1:]:
    rho_index2, theta_index2 = np.unravel_index(i, vote.shape)
    if abs(thetas[theta_index1] - thetas[theta_index2]) < np.deg2rad(5):
        break

rho1 = rhos[rho_index1]
theta1 = thetas[theta_index1]
rho2 = rhos[rho_index2]
theta2 = thetas[theta_index2]

medianrho   = (rho1 + rho2) / 2
mediantheta = (theta1 + theta2) / 2
\end{lstlisting}

This correctly identifies the two long edges of the tool (which are parallel) 
and averages their $\rho$ and $\theta$ to place the medial axis exactly between them.

\section{Final Approach}

\subsection{Stage 1: Background Subtraction}
\begin{lstlisting}
bgsub = cv.createBackgroundSubtractorMOG2(detectShadows=False)
subfr = bgsub.apply(frame)
\end{lstlisting}

MOG2 with \texttt{detectShadows=False} isolates the moving tool as a white 
foreground mask. Shadow detection is disabled as it adds unnecessary grey 
regions that complicate subsequent edge detection.

\subsection{Stage 2: Image Cleaning}
\begin{lstlisting}
medianbl = cv.medianBlur(subfr, 5)
gauss = cv.GaussianBlur(medianbl, (3, 3), 0.5)
close = cv.getStructuringElement(cv.MORPH_RECT, (3,3))
morph = cv.morphologyEx(gauss, cv.MORPH_CLOSE, close)
\end{lstlisting}

Median filtering first removes salt-and-pepper noise from the MOG2 mask. 
Gaussian blur then smooths the mask to improve Sobel gradient quality. 
Morphological closing fills small holes in the tool silhouette — important 
because thin tools can appear broken after background subtraction.

\subsection{Stage 3: Sobel Edge Detection}
\begin{lstlisting}
sx    = cv.Sobel(morph, cv.CV_64F, 1, 0, ksize=3)
sy    = cv.Sobel(morph, cv.CV_64F, 0, 1, ksize=3)
totder = np.sqrt(sx**2 + sy**2)
binary     = (totder > 50).astype(np.uint8)
\end{lstlisting}

Sobel operators compute the gradient magnitude at each pixel. The threshold 
of 50 on the magnitude selects only strong edges, reducing the number of 
edge pixels that must be voted in the Hough accumulator and improving 
peak sharpness.

\subsection{Stage 4: Custom Hough Line Transform}
For each edge pixel $(x, y)$ 
and each candidate angle $\theta$:
\begin{equation}
\rho = x\cos\theta + y\sin\theta
\end{equation}
The corresponding voting array cell $(\rho + d, \theta_{index})$ is incremented, 
where $d$ is the image diagonal used to offset $\rho$ to a non-negative index.

\begin{lstlisting}
height, width = binary.shape
diag = int(np.sqrt(height**2 + width**2))

rhos = np.arange(-diag, diag + 1, 1)
thetas = np.deg2rad(np.arange(0, 180))

vote = np.zeros((len(rhos), len(thetas)), dtype=np.uint64)
non_zero_y, non_zero_x = np.nonzero(binary)

for i in range(len(non_zero_x)):
    x = non_zero_x[i]
    y = non_zero_y[i]
    for theta_index in range(len(thetas)):
        theta = thetas[theta_index]
        rho = int(x * np.cos(theta) + y *np.sin(theta))
        rho_index = rho + diag
        vote[rho_index, theta_index] += 1
\end{lstlisting}

\subsection{Stage 5: Parallel Line Extraction and Medial Axis}
The voting array is sorted to find the strongest peak (first edge line), 
then the second-strongest peak within $5°$ degrees of the first (second parallel edge):

\begin{lstlisting}
index = np.argsort(vote.flatten())[::-1]
rho_index1, theta_index1 = np.unravel_index(index[0], vote.shape)
    
for i in index[1:]:
    rho_index2, theta_index2 = np.unravel_index(i, vote.shape)
    if abs(thetas[theta_index1] - thetas[theta_index2]) < np.deg2rad(5):
        break

rho1 = rhos[rho_index1],theta1 = thetas[theta_index1]
rho2 = rhos[rho_index2],theta2 = thetas[theta_index2]

medianrho = (rho1 + rho2) / 2
mediantheta = (theta1 + theta2) / 2

\end{lstlisting}

The medial axis line is drawn using the normal-form to Cartesian conversion:

\begin{align}
x_0 &= \cos(\theta_{med}) \cdot \rho_{med} \\
y_0 &= \sin(\theta_{med}) \cdot \rho_{med} \\
\Delta x &= -\sin(\theta_{med}), \quad \Delta y = \cos(\theta_{med})
\end{align}

\begin{lstlisting}
a = np.cos(mediantheta), b = np.sin(mediantheta)

x0 = a * medianrho, y0 = b * medianrho

dx = -b, dy = a

x1 = int(x0 + 1000 * dx), y1 = int(y0 + 1000 * dy)
x2 = int(x0 - 1000 * dx) ,y2 = int(y0 - 1000 * dy)

cv.line(frame, (x1, y1), (x2, y2), (255, 0, 255), 2)
\end{lstlisting}



\section{Results and Observations}
\begin{itemize}
    \item The medial axis line is correctly centered between the two detected 
    edges of the surgical tool in most frames.
    \item The $5°$ parallelism tolerance for the second line search successfully 
    rejects background edges (which occur at very different angles) while 
    accepting the second tool edge.
    \item Occasional frame-to-frame flicker occurs when the tool moves rapidly 
    and MOG2 produces an incomplete silhouette. 
    \item The nested loop implementation is computationally intensive. For real-time 
    use, numpy vectorisation of the accumulator fill loop (without changing 
    the algorithmic logic) would give a significant speedup.
\end{itemize}

\section{Future Work}
\begin{itemize}
    \item \textbf{Temporal smoothing:} An exponential moving average on the 
    detected $(\rho, \theta)$ parameters across frames would stabilise the 
    medial axis against frame-to-frame noise.
    \item \textbf{Probabilistic Hough:} A from-scratch implementation of the 
    probabilistic variant (random subset of edge pixels voted each frame) 
    would reduce computation while maintaining detection quality.
    \item \textbf{Adaptive edge threshold:} Replacing the fixed Sobel threshold 
    of 50 with Otsu's method on the gradient magnitude histogram would adapt 
    to lighting changes across the video.
    \item \textbf{MOG2 warmup:} Running 30 frames of background learning before 
    processing would provide a cleaner initial foreground mask.
\end{itemize}

\section{Conclusion}
A complete from-scratch Hough Line Transform pipeline was implemented for 
medial axis detection of a moving surgical tool. The algorithm correctly 
implements the $\rho = x\cos\theta + y\sin\theta$ voting, identifies 
the two dominant parallel tool edges, and computes their geometric average 
as the medial axis. Three non-trivial bugs were identified and fixed during 
development — most critically, the confusion between angle index and angle 
in radians when evaluating trigonometric functions, which caused the detected 
line to appear perpendicular to the tool. The pipeline satisfies the task 
constraint of zero built-in Hough transform usage.

\begin{thebibliography}{1}
\bibitem{numpy}
https://numpy.org/devdocs/

\bibitem{opencv_freecodecamp}
Beau Carnes, https://www.freecodecamp.org/news/opencv-full-course/

\bibitem{opencv_docs}
https://docs.opencv.org/

\bibitem{geeksforgeeks}
https://www.geeksforgeeks.org/

\bibitem{GitHubCommands}
Adam Lusted (etoxin), https://gist.github.com/etoxin/1acb257550b1de60fe99

\end{thebibliography}

\end{document}